\scanpage{20}
Les \unsure{choix|classes} de chemins de $\unsure{\xi_0|x_0}$ à $\unsure{\xi|x}$
définissent un torseur
\[
(3)\qquad P_{X,x}\quad\text{torseur à droite sous }\mathcal M_{g,\nu}
\quad\begin{gathered}(\text{\unsure{commutant}}\ \mathcal M_{X,x})\end{gathered}
\]
de telle façon que (1) soit déduit du diagramme analogue pour $(X_0,x_0)$
\[
(1_{g,\nu})\qquad
\diagram{\begin{tikzcd}
&1\arrow[d]&1\arrow[d]&&\\
&\widehat\Pi_{g,\nu}\arrow[r,equal]\arrow[d]&\widehat\Pi_{g,\nu}\arrow[d]&&\\
1\arrow[r,no head]&\widehat{\mathcal G}_{g,\nu}\arrow[r,no head]\arrow[d]&\mathcal M_{g,\nu}\arrow[r,no head]\arrow[d]&\Gamma_{\overline{\mathbb Q}/\mathbb Q}\arrow[r,no head]\arrow[d,equal]&1\\
1\arrow[r,no head]&\widehat\Gamma_{g,\nu}\arrow[r,no head]\arrow[d]&\mathcal N_{g,\nu}\arrow[r,no head]\arrow[d]&\Gamma_{\overline{\mathbb Q}/\mathbb Q}\arrow[r,no head]&1\\
&1&1&&
\end{tikzcd}}
\]
en le tordant, grâce à l'action de $\mathcal M_{g,\nu}$ opérant par autom. int.
sur lui-même, et sur sa filtration par $\widehat{\mathcal G}_{g,\nu},\widehat\Pi_{g,\nu}$.
En particulier
\[
(4)\quad\left\{\begin{aligned}
\Pi_{X,x}&\simeq P_{X,x}\wedge_{\mathcal M_{g,\nu}}\widehat\Pi_{g,\nu}\\
\mathcal G_{X,x}&\simeq P_{X,x}\wedge_{\mathcal M_{g,\nu}}\widehat{\mathcal G}_{g,\nu}
\qquad\deleted{\mathcal M_{X,x}}\\
\Gamma_{X,x}&\simeq P_{X,x}\wedge_{\mathcal M_{g,\nu}}\widehat\Gamma_{g,\nu}
\simeq\unsure{Q_X|Q_{X_0}}\wedge_{\mathcal N_{g,\nu}}\widehat\Gamma_{g,\nu}\\
\mathcal M_{X,x}&\simeq P_{X,x}\wedge_{\mathcal M_{g,\nu}}\mathcal M_{g,\nu}\\
\mathcal N_X&\simeq P_{X,x}\wedge_{\mathcal M_{g,\nu}}\mathcal N_{g,\nu}
\simeq Q_X\wedge_{\mathcal N_{g,\nu}}\mathcal N_{g,\nu}\\
\widetilde{\mathbb Q}&\simeq P_{X,x}\wedge_{\mathcal M_{g,\nu}}\overline{\mathbb Q}
\simeq Q_X\wedge_{\mathcal N_{g,\nu}}\overline{\mathbb Q}
\simeq R\wedge_{\Gamma_{\overline{\mathbb Q}/\mathbb Q}}\overline{\mathbb Q}
\end{aligned}\right.
\]
\mspage{19}
